% !PRINTED none
\begin{center}
\gesperrt{Zwanzigster Abschnitt.}

\gesperrt{Quadratische Körper.}
\end{center}

\begin{center}
\rule{0.2\linewidth}{0.4pt}
\end{center}

\begin{center}
\S.\,164.

\gesperrt{Basis eines quadratischen Körpers.}
\end{center}

Wir wollen die Resultate der allgemeinen Theorie der algebraischen Zahlen auf den besonderen Fall der Körper vom
$2^{\text{ten}}$ Grade, die auch quadratische Körper heissen, anwenden.
Ein quadratischer Körper $\Omega$ entspringt aus einer ganzzahligen
quadratischen Gleichung, und wird also aus dem Körper der
rationalen Zahlen durch Adjunction einer Quadratwurzel $\sqrt{d}$ abgeleitet.
Hierin kann $d$ als positive oder negative ganze Zahl
ohne quadratischen Theiler angenommen werden, und der einzige
 Werth $d = 1$ ist auszunehmen. Je nachdem $d$ positiv oder
negativ ist, erhalten wir einen reellen oder imaginären quadratischen
 Körper.

Den conjugirten Körper erhalten wir, wenn wir in allen
Zahlen in $\Omega$ $\sqrt{d}$ in $-\sqrt{d}$ verwandeln. Dadurch bekommen wir
aber keinen neuen Körper. Der quadratische Körper ist also
ein (absoluter) Normalkörper (\S.\,160).

Die Zahlen des Körpers $\Omega$ sind keine anderen als die in
\S.\,122 des ersten Bandes definirten quadratischen Irrationalzahlen,
die alle in der Form $x + y\sqrt{d}$ enthalten sind, worin
$x$ und $y$ ganze oder gebrochene rationale Zahlen bedeuten. Jede
irrationale Zahl $\omega$ dieses Körpers genügt einer quadratischen
Gleichung

\weq{1}{c\,\omega^2 = a + b\,\omega,}

in der $a, b, c$ ganze rationale Zahlen ohne gemeinsamen Theiler
sind, von denen $c$ als positiv angenommen werden kann, und
die Verbindung

\weq{2}{D = b^2 + 4\,a\,c}
%CONT>>
